\section[L-参数的构造]{L-参数的构造（2 月 12 日）}

\srcpage{143}

\subsection{设定}

设 $E$ 是非阿基米德局部域，$G/E$ 是约化群。固定一个素数 $\ell\neq p$。设 $\widehat G/\mathbb Z_\ell$ 是对偶群。正如我们讨论过的，$W_E$ 在 $\widehat G$ 上的作用有两种正规化。我们通过固定一个 $\sqrt q$ 的选取、并假装它落在 $\mathbb Z_\ell$ 中来回避这个问题。

\subsection{局部 Langlands 对应}

\subsubsection{表示论一侧}

我们关心光滑 $G(E)$-表示构成的范畴 $D(G(E),\mathbb Z_\ell)$。它全忠实地嵌入 $D_{\mathrm{lis}}(\Bun_G,\mathbb Z_\ell)$——后者是 $D_{\text{ét}}$ 的一个变体，对一切 $\mathbb Z_\ell$-代数 $\Lambda$ 都有效。（这里用到凝聚态数学的一些要素，特别是「固体六函子形式化」。）

\subsubsection{Galois 一侧}

我们考虑 L-参数的模叠 $\cZ^1(W_E,\widehat G)/\widehat G$（一个 Artin 叠）。

局部 Langlands 对应（local Langlands correspondence，下文简称 LLC）关注的正是这两侧之间的关系。在更经典的表述中，人们考察表示论一侧的不可约对象，把它们与 $\cZ^1(W_E,\widehat G)/\widehat G$ 的点相匹配。

\subsection{Bernstein 中心}

人们想要表达性质
\[
\text{「}\pi\mapsto\varphi_\pi\text{ 代数地变化」}.
\]

\begin{definition}
Bernstein 中心 $Z(G)$ 是光滑 $G(E)$-表示范畴上恒等函子的自同态构成的（交换）代数。
\end{definition}

具体地说，这意味着 $f\in Z(G)$ 对每个 $\pi$ 给出一个自同态 $f(\pi)\colon\pi\to\pi$，并且它与一切映射 $\pi\to\pi'$ 交换。

特别地，若 $f\in Z(G)$ 且 $\pi\in\operatorname{Irr}_{\overline{\mathbb Q}_\ell}(G)$，我们就得到标量 $f(\pi)\in\overline{\mathbb Q}_\ell$。这给出映射
\[
Z(G)_{\overline{\mathbb Q}_\ell}\hookrightarrow\{\operatorname{Irr}_{\overline{\mathbb Q}_\ell}(G)\text{ 上的函数}\},
\]
而 $Z(G)_{\overline{\mathbb Q}_\ell}$ 应当被看成「集合 $\operatorname{Irr}_{\overline{\mathbb Q}_\ell}(G)$ 上的代数函数」。

表达上述性质的一种方式是：我们希望对任意 $f\in\mathcal O(\cZ^1(W_E,\widehat G))^{\widehat G}$，映射 $\pi\mapsto f(\varphi_\pi)$ 都是「代数的」，即落在 Bernstein 中心 $Z(G)$ 中。

\begin{definition}
谱 Bernstein 中心是
\[
Z^{\mathrm{spec}}(G):=\mathcal O(\cZ^1(W_E,\widehat G))^{\widehat G}.
\]
\end{definition}

我们还考虑 $Z^{\mathrm{geom}}(G)$，即「$D_{\mathrm{lis}}(\Bun_G,\mathbb Z_\ell)$ 的 Bernstein 中心」，也就是 $\End(\operatorname{Id}_{D_{\mathrm{lis}}(\Bun_G,\mathbb Z_\ell)})$。把作用限制到 $G$ 的光滑表示上，我们得到映射
\[
Z^{\mathrm{geom}}(G)\to Z(G).
\]

\subsection{半单化的对应}

\begin{theorem}[Fargues--S]
若 $\ell$ 是 $\widehat G$ 的好素数，则存在典范映射
\[
\psi\colon Z^{\mathrm{spec}}(G)\to Z^{\mathrm{geom}}(G),
\]
它定义在 $\mathbb Z_\ell$ 上。（在 $\mathbb Q_\ell$ 上，不需要对 $\ell$ 作任何假设。）

特别地，若 $L/\mathbb Z_\ell$ 是代数闭域，$A\in D_{\mathrm{lis}}(\Bun_G,L)$ 且 $\End(A)=L$，则存在同态
\[
\varphi_A\colon W_E\to\widehat G(L),
\]
它在 $\widehat G(L)$-共轭的意义下唯一，是「半单的」，并且对一切 $f\in Z^{\mathrm{spec}}(G)$ 满足
\[
f(\varphi_A)=\psi(f)A\in L.
\]
\end{theorem}

\srcpage{144}

\begin{example}
若 $\pi$ 是 $G(E)$ 的不可约表示，则可取 $A=j_!\pi$。
\end{example}

\begin{remark}
我们只得到「半单的」L-参数，因为函数只能探测闭轨道。然而一般来说，人们会期望出现并非半单的 L-参数。
\end{remark}

\begin{proposition}[对应的性质]
映射 $\pi\mapsto\varphi_\pi$ 具有以下性质。
\begin{enumerate}[label=(\arabic*)]
\item 对环面，它与通常的 LLC 一致。
\item 它与扭转和中心特征标相容。
\item 它与对偶相容。
\item 若 $G'\to G$ 是诱导出伴随群同构的映射，$\pi$ 是 $G(E)$ 的不可约表示，$\pi'$ 是 $\pi|_{G'(E)}$ 的不可约成分，则 $\varphi_{\pi'}$ 是 $\varphi_\pi$ 在 $\widehat G\to\widehat G'$ 下的像。
\item 它与乘积相容。
\item 它与纯量 Weil 限制相容。
\item 它与抛物诱导相容。
\item 对 $\mathrm{GL}_n$，它与通常的 LLC 一致。\footnote{这是唯一用到整体方法的地方；事实上，通常的 LLC 正是用整体方法证明的。}
\item 它与 $\Bun_G$ 上的 Hecke 函子相容。特别地，可以计算 $A$ 的 Hecke 修正的 L-参数。
\item 它与局部 shtuka 模空间的上同调相容，例如局部 Shimura 簇（特别是 Rapoport--Zink 空间）。
\end{enumerate}
\end{proposition}

\begin{corollary}[Helm--Moss 定理]
设 $G=\mathrm{GL}_n$，则由通常的 LLC 定义的映射
\[
Z^{\mathrm{spec}}(G)_{\mathbb Q_\ell}\to Z(G)_{\mathbb Q_\ell}
\]
实际上可整系数地定义，即诱导出 $Z^{\mathrm{spec}}(G)\to Z(G)$。这表达了「LLC 与同余的相容性」。
\end{corollary}

\subsection{$\psi\colon Z^{\mathrm{spec}}(G)\to Z^{\mathrm{geom}}(G)$ 的构造}

我们先总结手上已有的数据。

\begin{itemize}
\item 我们有 $\infty$-范畴 $\mathcal C=D_{\mathrm{lis}}(\Bun_G,\mathbb Z_\ell)$，并且对任意有限集 $I$，有精确幺半函子（选取一个有限商 $W_E\twoheadrightarrow Q\curvearrowright\widehat G$，使得对 $\widehat G$ 的作用通过它分解）
\[
\Rep_{\mathbb Z_\ell}(\widehat G\rtimes Q)^I\to\End(\mathcal C)^{W_E^I},
\]
记作 $V\mapsto T_V$；它是 $\Rep_{\mathbb Z_\ell}(Q^I)$-线性的，并且对 $I$ 具有函子性。这来自 Hecke 作用。
\end{itemize}

我们只会用到这类抽象数据。

\begin{proposition}
对任意 $A\in\mathcal C^\omega$，存在野惯性子群 $I_E$ 的开子群 $P$，使得对一切 $I$ 与 $V\in\Rep(\widehat G\rtimes Q)^I$，$T_V(A)$ 上的 $W_E^I$-作用都通过 $(W_E/P)^I$ 分解。
\end{proposition}

这意味着我们在上面可以把 $W_E$ 换成 $W_E/P$。于是（如上次那样）可以把 $W_E/P$ 换成一个离散化 $W\subset W_E/P$。

上次我们陈述了：

\srcpage{145}

\begin{theorem}
在温和的技术条件下，映射
\[
\varinjlim_{(n,F_n\to W)}\mathcal O(\widehat G^n)^{\widehat G}\to\mathcal O(\cZ^1(W_E/P,\widehat G))^{\widehat G}
\]
是同构。
\end{theorem}

我们想找到映射 $\mathcal O(\cZ^1(W_E/P,\widehat G))^{\widehat G}\to Z^{\mathrm{geom}}(G)=\End(\operatorname{Id}_{\mathcal C})$。由上述定理，只需给出映射 $\varinjlim_{(n,F_n\to W)}\mathcal O(\widehat G^n)^{\widehat G}\to Z^{\mathrm{geom}}(G)$。这将通过「巡游算子」来实现（V. Lafforgue）。

\begin{definition}
(1) 一个巡游数据是元组 $(I,V,\alpha,\beta,(\gamma_i)_{i\in I})$，其中
\begin{itemize}
\item $I$ 是有限集；
\item $V\in\Rep(\widehat G\rtimes Q)^I$；
\item $\alpha\colon\mathbf 1\to V|_{\Delta(\widehat G)}$，$\beta\colon V|_{\Delta(\widehat G)}\to\mathbf 1$；
\item $\gamma_i\in W_E$。
\end{itemize}
\hspace*{2em}(2) 给定巡游数据，巡游算子是 $\End(\operatorname{Id}_{\mathcal C})$ 中的如下元素。对任意 $A\in\mathcal C$，诱导出的 $A$ 的自同态为
\[
A=T_{\mathbf 1}(A)\xrightarrow{\alpha}T_V(A)\xrightarrow{(\gamma_i)_{i\in I}}T_V(A)\xrightarrow{\beta}T_{\mathbf 1}(A)=A.
\]
\end{definition}

\begin{proposition}
巡游算子的全体定义了映射
\[
\varinjlim_{(n,F_n\to W)}\mathcal O(\widehat G^n)^{\widehat G}\to Z^{\mathrm{geom}}(G).
\]
\end{proposition}

\begin{corollary}
L-参数 $\varphi_A$ 的刻画如下：对一切巡游数据，标量
\[
L\xrightarrow{\alpha}V\xrightarrow{(\varphi_A(\gamma_i))_{i\in I}}V\xrightarrow{\beta}L
\]
与标量
\[
A=T_{\mathbf 1}(A)\xrightarrow{\alpha}T_V(A)\xrightarrow{(\gamma_i)_{i\in I}}T_V(A)\xrightarrow{\beta}T_{\mathbf 1}(A)=A
\]
一致。
\end{corollary}

\subsection{谱作用}

\begin{theorem}
上述 $\infty$-范畴数据等价于 $\Perf(\cZ^1(W_E,\widehat G)/\widehat G)$ 在 $D_{\mathrm{lis}}(\Bun_G,\mathbb Z_\ell)$ 上的一个作用。
\end{theorem}

\begin{remark}
Nadler--Yun 以及 Gaitsgory--Kazhdan--Rozenblyum--Varshavsky 有相关工作。这里的一个新颖之处在于我们在整系数下工作。
\end{remark}

\begin{example}
我们来说明这如何恢复 Bernstein 中心的映射。存在函子
\[
\Rep(\widehat G\rtimes Q)^I\to\Perf(\cZ^1(W_E,\widehat G)/\widehat G)^{W_E^I},
\]
它由拉回与张量给出。定理 27.13 于是给出到 $\End(\mathcal C)^{W_E^I}$ 的函子。
\end{example}

\subsection{椭圆 L-参数}

对「椭圆」L-参数而言，这意味着什么？

为简单起见，设 $G$ 半单，系数取 $\overline{\mathbb Q}_\ell$。称 $\varphi$ 是椭圆的，若它定义了 $\cZ^1(W_E,\widehat G)_{\overline{\mathbb Q}_\ell}/\widehat G$ 的一个孤立分支；即：$S_\varphi\subset\widehat G$ 是它的中心化子，且包含
\[
[*\,/S_\varphi]\subset[\cZ^1(W_E,\widehat G)/\widehat G]
\]
是开且闭的。特别地，它由某个幂等元切出。因此由谱作用，我们得到相应的分支
\[
D^\varphi_{\mathrm{lis}}(\Bun_G,\overline{\mathbb Q}_\ell)\overset{\oplus}{\subset}D_{\mathrm{lis}}(\Bun_G,\overline{\mathbb Q}_\ell).
\]

\srcpage{146}

现在我们在 $D_{\mathrm{lis}}(\Bun_G,\overline{\mathbb Q}_\ell)$ 上有谱作用，而作用在 $D^\varphi_{\mathrm{lis}}(\Bun_G,\overline{\mathbb Q}_\ell)$ 上的是 $\Rep(S_\varphi)$。谱作用与 Hecke 作用的相容性说的是：给定 $V\in\Rep(\widehat G\rtimes Q)^I$，若分解为
\[
V|_{S_\varphi\times W_E}=\bigoplus_{i=1}^{r}W_i\boxtimes r_i,\qquad W_i\in\operatorname{Irr}\Rep(S_\varphi),\quad r_i\in\Rep_{\overline{\mathbb Q}_\ell}(W_E),
\]
则对 $A\in D^\varphi_{\mathrm{lis}}(\Bun_G,\overline{\mathbb Q}_\ell)$ 有
\[
T_V(A)=\bigoplus_{i=1}^{m}\operatorname{Act}_{W_i}(A)\otimes r_i\in\left(D^\varphi_{\mathrm{lis}}(\Bun_G,\overline{\mathbb Q}_\ell)\right)^{W_E},
\]
其中 $\operatorname{Act}_{W_i}(A)$ 指 $\Rep(S_\varphi)$ 在 $D^\varphi_{\mathrm{lis}}(\Bun_G,\overline{\mathbb Q}_\ell)$ 上的作用。

在某种意义上，这「就是」Kottwitz 猜想。

\begin{proposition}
所有 $A\in D^\varphi_{\mathrm{lis}}(\Bun_G,\overline{\mathbb Q}_\ell)$ 都集中在半稳定轨迹上，并对应超尖点表示。
\end{proposition}

\begin{proof}[证明]
这由与抛物诱导的相容性推出：非半稳定的各层与 Levi 子群相关，而后者的参数必须来自抛物诱导。
\end{proof}

我们有分解
\[
D^\varphi_{\mathrm{lis}}=\bigoplus_{b\in B(G)_{\mathrm{basic}}}\bigoplus_{\substack{\pi\ \text{不可约超尖点}\\\text{$G_b(E)$ 的表示}\\\varphi_\pi=\varphi}}D(\overline{\mathbb Q}_\ell)\cdot[\pi].
\]

它有 $\Rep(S_\varphi)$ 的作用。若 $S_\varphi$ 交换（经典类型中就是如此），$S_\varphi$ 的特征标将置换这些 $\pi$。这是一个「Jacquet--Langlands 对应」。

此外，$T_V([\pi])$ 是某个局部 shtuka 模空间上同调中的 $\pi$-同型分量，并且
\[
T_V([\pi])=\bigoplus_{i=1}^{m}\operatorname{Act}_{W_i}([\pi])\otimes r_i,
\]
其中 $\operatorname{Act}_{W_i}([\pi])$ 是 Jacquet--Langlands 转移。这与 Kottwitz 猜想的预测相符。然而，我们并未证明 Kottwitz 猜想，因为我们还缺少对一切满足 $\varphi_\pi=\varphi$ 的 $\pi$ 的参数化。

\begin{conjecture}
设 $G$ 拟分裂。固定 Whittaker 数据。则存在唯一的一般表示 $\pi_0$，使 $\varphi_{\pi_0}=\varphi$，且函子
\[
\Perf([*\,/S_\varphi])\to(D^\varphi_{\mathrm{lis}})^\omega=\bigoplus_b\bigoplus_\pi\Perf(\overline{\mathbb Q}_\ell)[\pi],
\]
$W\mapsto\operatorname{Act}_W([\pi_0])$，是范畴等价。
\end{conjecture}

特别地，这诱导出双射
\[
\operatorname{Irr}\Rep(S_\varphi)\xrightarrow{\sim}\{\pi\}.
\]

\begin{remark}
这是 Kaletha 利用 $B(G)_{\mathrm{basic}}$ 对 LLC 的表述。
\end{remark}

\srcpage{147}

\subsection{回到 $D_{\mathrm{lis}}$}

我们现在回到关于 $\Bun_G$ 上层的整个范畴的一个猜想。

固定 Whittaker 数据，即 $U\subset B\subset G$，以及一个非退化特征标 $\psi\colon U(E)\to\breve{\mathbb Z}_\ell^\times$。于是我们有表示 $c\text{-}\mathrm{Ind}_{U(E)}^{G(E)}(\psi)$。

\begin{conjecture}
存在等价
\[
D^b_{\mathrm{coh}}(\cZ^1(W_E,\widehat G)_{\overline{\mathbb Q}_\ell}/\widehat G)\cong D_{\mathrm{lis}}(\Bun_G,\overline{\mathbb Q}_\ell)^\omega,
\]
它在 $\Perf(\cZ^1(W_E,\widehat G)_{\overline{\mathbb Q}_\ell}/\widehat G)$ 上线性，并把 $\mathcal O_{\cZ^1(W_E,\widehat G)_{\overline{\mathbb Q}_\ell}/\widehat G}$ 送到 $[c\text{-}\mathrm{Ind}_{U(E)}^{G(E)}(\psi)]$。
\end{conjecture}

\begin{remark}
谱作用决定了猜想 27.19 的函子在子范畴
\[
\Perf(\cZ^1(W_E,\widehat G)_{\overline{\mathbb Q}_\ell}/\widehat G)\subset D^b_{\mathrm{coh}}(\cZ^1(W_E,\widehat G)_{\overline{\mathbb Q}_\ell}/\widehat G)
\]
上的取值。$D^b_{\mathrm{coh}}(\cZ^1(W_E,\widehat G)_{\overline{\mathbb Q}_\ell}/\widehat G)$ 与 $\Perf(\cZ^1(W_E,\widehat G)_{\overline{\mathbb Q}_\ell}/\widehat G)$ 的差别来自 $\cZ^1(W_E,\widehat G)_{\overline{\mathbb Q}_\ell}$ 的奇点。
\end{remark}

\begin{remark}
在整系数层面，我们期望存在范畴等价，但需要在 Langlands 参数叠上施加「幂零奇异支撑」条件。我们期望的等价是
\[
D^b_{\mathrm{coh},\mathrm{Nilp}}(\cZ^1(W_E,\widehat G)_{\breve{\mathbb Z}_\ell}/\widehat G)\cong D_{\mathrm{lis}}(\Bun_G,\breve{\mathbb Z}_\ell)^\omega.
\]
\end{remark}
